\section{磁化を符号・秩序・履歴に分ける}\label{theory:em:magnetism}
単位体積あたりの磁気モーメントを磁化\(\mathbf M\)とする。SIでは
\(\mathbf B=\mu_0(\mathbf H+\mathbf M)\)。等方的な弱い線形応答では
\(\mathbf M=\chi\mathbf H\)で、反磁性は負の\(\chi\)、常磁性は小さな正の\(\chi\)を持つ。
強磁性は外場なしの自発磁化を持つ秩序であり、多磁区試料の正味磁化が小さいことと両立する。
反強磁性は典型的には反平行の副格子磁化が打ち消し合う秩序で、Néel温度を超えると秩序を失う。
通常の強磁性の秩序を失う温度はCurie温度と呼ぶ。残留磁化は外場を取り去った後の履歴による磁化で、
磁区配置と保磁力を含む材料の状態に依存する。\cite{openstax_magnetism}

室温の代表的強磁性元素はFe・Co・Niである。過去問のMnをCoに読み替えない。
磁性は温度・結晶相・組成にも依存する。中性子回折は核による構造散乱だけでなく磁気散乱にも使える。
Slater--Pauling曲線は合金の電子数と磁気モーメント、Ising模型はスピン相互作用と秩序を扱う。
一方、Wiedemann--Franz則の\(\kappa_e/(\sigma T)\)は電子の熱電輸送係数の比である。
磁性に関する相対比較で区別し、磁場やスピンの影響を全く受けないと断定しない。\cite{em51_mit_magnetic_readings,em51_wf}

\section{局在モーメントと伝導電子}\label{theory:em:pauli}
独立した局在モーメント、弱い外場、高温側のモデルでは、磁化の方向が熱揺らぎで乱され、
磁化率はCurie則\(\chi=C/T\)に従う。Tは絶対温度である。
相互作用を含むCurie--Weiss近似の\(C/(T-\theta)\)も、転移点近傍や全温度の厳密則ではない。
「常磁性なら必ず同じ温度依存」という暗記を避ける。

Pauli常磁性は金属の伝導電子のスピンへの応答で、占有のFermi統計が重要になる。
低温では全電子が独立に方向を変える局在モーメントのモデルを使えず、Fermi準位近傍の状態が応答する。
原子核まわりの軌道運動を起源とするという過去問の説明はこの機構と違う。
自由電子の軌道応答であるLandau反磁性、束縛電子の応答を分ける。\cite{em51_pauli}

超伝導のMeissner状態では磁束が排除されるが、温度・磁場の条件がある。
II型の混合状態では量子化磁束が侵入するため、「超伝導の全状態で完全排除」とはしない。
電気抵抗ゼロという輸送性質だけから、磁化の全状態を決めない。

\section{形状で変わる反磁場と磁場勾配の力}\label{theory:em:demag}
試料の磁化は周囲と内部の場を変える。均一磁化した楕円体では、主軸方向ごとに
\[
 H_{\mathrm d,i}=-N_iM_i,\qquad N_x+N_y+N_z=1
\]
と書ける。NはSIでの無次元の反磁場係数で、形状に依存する。
一般の不均一磁化・任意形状を一つのスカラーNで表すとは限らない。\cite{em51_demag}
反磁場は磁化から生じる場であり、反磁性という材料応答の名前とは別である。

黒鉛にも反磁性があるため、磁場中で「全く力を受けない」という全称は誤りである。
弱い等方的な誘起磁化の小試料なら、背景場を変えない近似で
\(\mathbf F\simeq\chi V\nabla(B^2)/(2\mu_0)\)と考えられる。
反磁性では強い場から弱い場へ向かう。実際の黒鉛の異方性や試料形状は別に考える。
理想的な一様場の並進力ゼロと、非一様場で働き得る力を区別する。

\section{巻数比を磁束変化から導く}\label{theory:em:transformer}
理想変圧器の一次・二次が共通の磁束\(\Phi(t)\)を共有するなら、
誘導起電力の大きさはそれぞれ\(N_1|d\Phi/dt|\)、\(N_2|d\Phi/dt|\)である。
したがって同じ位相規約での電圧の大きさ、実効値又は最大値同士では
\[
 \frac{E_2}{E_1}=\frac{N_2}{N_1}.
\]
巻数が多い側ほど電圧が大きいことが検算となる。損失を無視した電力保存から
\(E_1I_1=E_2I_2\)なので電流比は逆になる。
磁束漏れ、巻線抵抗、飽和、周波数応答を含む実機の限界は、この理想比の条件の外である。\cite{em51_transformer_ref}

\section{誘導加熱と誘電加熱を損失で分ける}\label{theory:em:heating}
誘導加熱は時間変化する磁場で導電体に渦電流を作り、そのJoule損を利用する。
磁化履歴損は強磁性体等に寄与するが、すべての導電体に必須の損失ではない。
誘電加熱は交流電場に対する分極応答の損失を利用し、代表的な対象は誘電体である。
有限の導電率を持つ材料では導電損と分極損が共存する場合がある。
マイクロ波加熱も誘電損失を利用できるが、高周波加熱の全方式がGHz帯という意味ではない。\cite{em51_heating}

良導体の線形な半空間に入る調和場では、表面からの場・電流の振幅が
\(\exp(-z/\delta)\)の形で減衰する。代表的な表皮深さは
\[
 \delta=\sqrt{\frac{2}{\omega\mu\sigma}}
       =\sqrt{\frac{\rho}{\pi f\mu}} .
\]
深さ\(\delta\)で振幅が\(1/e\)になるため、内部への侵入が厳密にゼロではない。
抵抗率・透磁率・周波数の条件を読み、有限試料の全温度分布をこの式だけで決めない。
加熱された表面から内部への熱伝導は電磁場の侵入とは別の過程である。\cite{em51_heating,em51_skin}
プラズマも高周波場から加熱できる。融合プラズマのRF加熱、無電極放電トーチ、
固体の渦電流加熱の回路と粒子の吸収機構は同一ではない。\cite{em51_plasma}

\section{副殻の容量と半導体の四面体結合}\label{theory:em:electrons}
一個の空間軌道にはスピンが異なる最大2電子を収容する。
s・p・d・f副殻の空間軌道数は1・3・5・7で、総容量は2・6・10・14である。
原典の「軌道」の語を、一個の軌道と副殻群で区別する。
Ge原子の外殻配置\(4s^24p^2\)と、通常の結晶の4近接・四面体のsp3結合モデルを分ける。
sp2の三角平面構造をダイヤモンド型へ当てはめない。\cite{atomic_quantum,em51_ge,em51_ge_structure}

SiへBをドープするとアクセプターを介して正孔が主要な担い手になる。
金属での負電子の平均ドリフト、半導体の正孔、孤立原子のスピンと
バルク金属のPauli応答を区別する。\cite{mit_semiconductor,em51_pauli}

\section{圧電性の逆命題を使わない}\label{theory:em:ferroelectric}
圧電性は応力による電気応答と電場によるひずみの結合である。
強誘電性は外場なしの自発分極と、その向きの電場による反転を特徴とする。
石英は圧電性を持つが、通常の強誘電体のような反転する自発分極を持たない。
圧電性から強誘電性を保証する逆命題は成立しない。\cite{mit_ferroelectric,em51_quartz}
多ドメイン試料では局所の自発分極が存在しても、正味の分極は相殺され得る。

典型的な強誘電体ではCurie温度を越えて常誘電相へ移る。
一次相転移は秩序変数が不連続に変わる場合で、強誘電相転移では自発分極を用いて考える。
有限の測定電場での飽和値・履歴と、電場ゼロでの熱力学的な秩序変数を無条件に同一視しない。
原子の変位、配向の秩序、結晶対称性と分極をつなぎ、微視的な機構を全物質で一種類に固定しない。\cite{em51_ferro_transition}

\section{特性線・連続放射と法令上の電波}\label{theory:em:spectral}
真空中の電磁波では\(c=f\lambda\)、光子エネルギーは\(E=hf\)である。
赤外・可視・紫外の順に波長が短くなる。
X線管の放射は、電子の減速による連続的な制動放射に、内殻準位差に対応する特性線が重なる。
管全体のスペクトルが連続成分を含むことと、特性X線自体を連続スペクトルと呼ぶことは違う。\cite{em51_xray}
取得できた電波法第二条の周波数上限は300万MHz以下で、\(3\times10^{12}\)Hz以下へ換算される。
令和5年度の施行版の直接取得には失敗したため、その時点の条文照合は個別問題で未確定として残す。
用途に基づくマイクロ波の波長区分と、法令の電波の範囲を別々に読む。\cite{em51_radio_law}
